\chapter{Metode Penelitian}
\emph{Cara menulis bab ini.} Bab Metode Penelitian (Sistem Informasi) menjelaskan data/kebutuhan, metodologi pengembangan sistem, serta alat dan bahan. Uraikan cukup rinci agar dapat direplikasi.


\section{Data dan Kebutuhan Sistem}
\emph{Cara menulis.} Uraikan kebutuhan fungsional dan non-fungsional sistem serta data yang digunakan.

Bagian ini menjelaskan sumber data (mis.\ basis data BPS) serta kebutuhan sistem
yang akan dikembangkan.

\section{Metodologi Pengembangan Sistem}
\emph{Cara menulis.} Jelaskan metodologi yang dipilih (mis.\ \emph{Waterfall}, \emph{Prototyping}, \emph{Agile}) beserta tahapannya.

Pengembangan sistem mengikuti pendekatan \emph{System Development Life Cycle}
(SDLC) dengan model \emph{prototyping}. Tahapan pengembangan disajikan pada
Gambar~\ref{gbr:sdlc}.

\begin{figure}[H]
\centering
\begin{tikzpicture}[node distance=0.95cm]
  \node[fc-mulai] (a) {Komunikasi/Analisis Kebutuhan};
  \node[fc-proses, below=of a] (b) {Perancangan Cepat (Prototipe)};
  \node[fc-proses, below=of b] (c) {Pembangunan Prototipe};
  \node[fc-putusan, below=of c] (d) {Diterima?};
  \node[fc-proses, below=of d] (e) {Implementasi \& Pengujian};
  \draw[fc-panah] (a)--(b); \draw[fc-panah] (b)--(c); \draw[fc-panah] (c)--(d);
  \draw[fc-panah] (d)-- node[right]{ya} (e);
  \draw[fc-panah] (d.west) -- ++(-2.2,0) |- node[pos=0.25,left]{tidak} (b.west);
\end{tikzpicture}
\caption{Tahapan pengembangan sistem dengan model prototyping}
\label{gbr:sdlc}
\end{figure}

\section{Alat dan Bahan}
\emph{Cara menulis.} Sebutkan perangkat keras/lunak, bahasa, dan \emph{framework} yang digunakan.

Perangkat lunak dan pustaka yang digunakan, antara lain bahasa pemrograman,
kerangka kerja (\emph{framework}), dan sistem manajemen basis data.
